\documentclass[10pt,a4paper]{amsart}
\usepackage[margin=19mm]{geometry}
\usepackage{amsmath,amssymb,mathtools}
\usepackage[scr=rsfs,cal=euler]{mathalfa}
\usepackage{xfrac}
\usepackage[unicode,hidelinks,bookmarks=false]{hyperref}
\newcommand{\ie}{i.e.\ }
\newcommand{\cf}{cf.\ }
\newcommand{\resp}{respectively\ }
\newcommand{\Subsection}{\subsection}
\providecommand{\nobreakdash}{\nobreak\hbox{-}}
\setlength{\emergencystretch}{2em}
\multlinegap0pt
\usepackage{amssymb}
\usepackage{tikz}
\newtheorem{lemma}{Lemma}
\title{Homological Mirror Symmetry for Affine Log Calabi-Yau Surfaces}
\author{Umut Varolg\"une\c{s}}
\date{Source: arXiv:2608.13779v1 (CC BY 4.0)}
\begin{document}
\maketitle

\noindent\emph{Source and licence.} Homological Mirror Symmetry for Affine Log Calabi-Yau Surfaces. Version arXiv:2608.13779v1 (CC BY 4.0), distributed by arXiv under the Creative Commons Attribution 4.0 International licence, http://creativecommons.org/licenses/by/4.0/, as recorded in the arXiv licence metadata for this exact version and shown on the abstract page https://arxiv.org/abs/2608.13779v1. Full licence notice: \texttt{licenses/arxiv-2608.13779-CC-BY-4.0.txt}.\par\smallskip

\noindent\textit{Editorial scope.} Counted unit: the article's lemma lemma (source0.tex lines 490--499). The article's complete original proof of it is delivered with it (source0.tex lines 501--517, one \texttt{proof} environment of 17 lines). No mathematics is written, repaired or re-derived: the statement and the proof are the article's own lines, in the article's own notation, and no retained line is substituted or re-flowed.  No further source-local context is needed: every cross-reference of every retained line resolves inside the two delivered spans.  Nothing was omitted from any retained span, so the package carries no omission record.  No retained line contains a citation, so the wrapper carries no bibliography and no bibliography entry.  No retained line contains a figure environment, an image-inclusion command or drawing code, so no image is delivered and the package declares no asset dependency.  Editorial formatting only, in addition: an AMS \texttt{amsart} preamble that restates the article's own declarations and macros used by the retained lines, namely the environments \texttt{lemma} on separate counters, together with the standard packages those lines need.  The retained environments print in this document as Lemma 1 in order of appearance, whereas the article prints the same items with the numbering of its own full document.  The complete native source of the article as distributed in the arXiv e-print for this exact version is bundled unchanged as \texttt{sources/207612/source0.tex}.
\begin{lemma}[Resolving Triangle Packing]
For a sufficiently small $\epsilon > 0$, we define the shifted big triangle for each edge by shifting the entire big triangle parallel to $\vec{v}_i$ by $\epsilon$ times the inverse of the affine distance of $p$ to $e_i$. The apex of the shifted big triangle becomes
\begin{equation}
    A_i := p + \frac{\epsilon}{d(p, e_i)} \vec{v}_i.
\end{equation} 

Here sufficiently small only means that, after the shifts, triangles are still contained in $\Delta_X$ and the base edges are contained in the interiors of the $e_i$. 

Then, the shifted big triangles are automatically disjoint.
\end{lemma}

\begin{proof}
We first notice that the vector connecting $A_{i-1}$ to $A_i$ is a positive scalar multiple of the vector pointing from $R_i$ to $p$. One way to see this is to apply an integral affine transformation that maps the corner $R_i$ to the origin $(0,0)$. We can locally align $e_i$ with the positive $x$-axis and $e_{i-1}$ with the positive $y$-axis (meaning $\vec{v}_i = (1,0)$ and $\vec{v}_{i-1} = (0,-1)$).

Let $p = (x,y)$. By definition, $d(p, e_{i-1}) = x$ and $d(p, e_i) = y$. Evaluating the shifts yields:
\begin{align*}
    A_{i-1} &= p + \frac{\epsilon}{x} (0,-1) = \left(x, y - \frac{\epsilon}{x}\right) \\
    A_i &= p + \frac{\epsilon}{y} (1,0) = \left(x + \frac{\epsilon}{y}, y\right)
\end{align*}
The vector connecting $A_{i-1}$ to $A_i$ is calculated as:
\begin{equation}
    A_i - A_{i-1} = \left(\frac{\epsilon}{y}, \frac{\epsilon}{x}\right) = \frac{\epsilon}{xy} (x,y) = \frac{\epsilon}{d(p, e_{i-1}) d(p, e_i)} (p - R_i),
\end{equation} proving the desired claim.

Because $\Delta_X$ is strictly convex, the vectors $p - R_i$ rotate counterclockwise by less than $\pi$ radians as we traverse the corners. Consequently, the edges of the sequence $A_1, \dots, A_m$ also rotate counterclockwise by less than $\pi$ radians, proving that they form a \emph{strictly convex} polygon $P_{\text{micro}}$.

The extended edges of any strictly convex polygon partition the complement into a set of ``exterior wedges.'' Here by exterior wedge we specifically take those wedges bounded by the rays obtained by extending each edge $A_iA_{i-1}$ in the direction of $A_{i-1}$. We now notice that the  $i$-th shifted triangle is completely confined within the $i$-th exterior wedge of $P_{\text{micro}}$ and, except for the vertex $A_i$, is contained in the interior of that exterior wedge. Because the unions of the interiors with their vertices (of each wedge) are pairwise disjoint, the proof is complete.
\end{proof}

\end{document}
